% © 2026 Ing. David Jaroš — CC BY-NC-ND 4.0
%
% This work is licensed under a Creative Commons Attribution-NonCommercial-NoDerivatives
% 4.0 International License (CC BY-NC-ND 4.0).
%
% N_eff Derivation from UBT Action S[Θ] — Step 2: One-Loop Vacuum Polarization
% Track: CORE — α derivation / B-coefficient
% Date: 2026-03-05

\ifdefined\INCLUDEMODE
  % Being included in another document — skip preamble
\else
  \documentclass[11pt,a4paper]{article}
  \usepackage{amsmath,amssymb,amsthm}
  \usepackage{hyperref}
  \usepackage{booktabs}

  \newcommand{\C}{\mathbb{C}}
  \newcommand{\R}{\mathbb{R}}
  \renewcommand{\H}{\mathbb{H}}
  \newcommand{\Tr}{\mathrm{Tr}}
  \newcommand{\Sc}{\mathrm{Sc}}

  \newtheorem{theorem}{Theorem}
  \newtheorem{lemma}[theorem]{Lemma}
  \newtheorem{proposition}[theorem]{Proposition}
  \newtheorem{definition}[theorem]{Definition}
  \newtheorem{corollary}[theorem]{Corollary}
  \theoremstyle{remark}
  \newtheorem{remark}{Remark}

  \title{N$_{\mathrm{eff}}$ Derivation from UBT Action $S[\Theta]$\\
    \large Step 2: One-Loop Vacuum Polarization from $S_{\mathrm{kin}}$}
  \author{Ing.\ David Jaroš}
  \date{2026-03-05}

  \begin{document}
  \maketitle
\fi

\section{Step 2: One-Loop Vacuum Polarization from $S_{\mathrm{kin}}$}
\label{sec:step2_vacuum_polarization}

\subsection{Goal}

Compute the one-loop photon vacuum polarization $\Pi(k^2)$ generated by integrating out the
biquaternionic field $\Theta$ in the kinetic action $S_{\mathrm{kin}}$, and identify the
coefficient $N_{\mathrm{eff}}$ in
\begin{equation}
  B_0 = \frac{2\pi\,N_{\mathrm{eff}}}{3}
  \label{eq:B0_Neff}
\end{equation}
as the number of independent charged complex d.o.f.\ of $\Theta$.

\subsection{Starting Action}

From Appendix~A (\texttt{appendix\_AA\_theta\_action.tex}) and the mode decomposition of
Step~1, the UBT kinetic action coupled to a background electromagnetic field $A_\mu$ is:
\begin{equation}
  S_{\mathrm{kin}}[\Theta, A]
  = \int d^4x\,d\psi\;\Sc\!\left[(D_\mu\Theta)^\dagger(D^\mu\Theta)\right],
  \label{eq:Skin_UBT}
\end{equation}
where $D_\mu = \partial_\mu + ie A_\mu$ is the covariant derivative and $\Sc[\cdot]$
extracts the quaternion-scalar part.

After expanding $\Theta = \sum_{a=0}^{3} z_a E_a$ in the $E_a$ basis (Step~1,
Sec.~\ref{sec:step1_mode_decomposition}), the kinetic action becomes a sum of
$N_{\mathrm{charged}}$ independent complex scalar kinetic terms:
\begin{equation}
  S_{\mathrm{kin}}[\{z_a\}, A]
  = \sum_{a=0}^{N_{\mathrm{charged}}-1}
    \int d^4x\,d\psi\;(D_\mu z_a)^*(D^\mu z_a),
  \label{eq:Skin_sum}
\end{equation}
where each $z_a$ is a complex scalar with charge $+e$.

\subsection{Expansion in the Background Field}

Expand $S_{\mathrm{kin}}$ to second order in $A_\mu$. For a single complex scalar $\phi$
with $D_\mu = \partial_\mu + ie A_\mu$:
\begin{equation}
  (D_\mu \phi)^*(D^\mu \phi)
  = (\partial_\mu \phi^*)(\partial^\mu \phi)
    + ie A_\mu[\phi^*(\partial^\mu \phi) - (\partial^\mu \phi^*)\phi]
    + e^2 A_\mu A^\mu |\phi|^2.
  \label{eq:Dmu_expand}
\end{equation}
The $O(A^2)$ term determines the vacuum polarization:
\begin{equation}
  S_{\mathrm{kin}}^{(2)}[A]
  = \frac{1}{2}\int \frac{d^4k}{(2\pi)^4}\;A_\mu(-k)\,\Pi^{\mu\nu}(k)\,A_\nu(k),
  \label{eq:S2_def}
\end{equation}
where the vacuum polarization tensor is transverse:
\begin{equation}
  \Pi^{\mu\nu}(k) = (k^2 g^{\mu\nu} - k^\mu k^\nu)\,\Pi(k^2).
  \label{eq:Pimunu_transverse}
\end{equation}

\subsection{One-Loop Result for a Single Complex Scalar}

For a single complex scalar field $\phi$ of charge $e$, mass $m$, the standard dimensional
regularization result is (see e.g.\ Peskin \& Schroeder, Chapter~7):
\begin{equation}
  \Pi^{(1)}(k^2)
  = \frac{e^2}{2\pi^2}\int_0^1 dx\; x(1-x)
    \left[\frac{2}{\epsilon} - \ln\!\left(\frac{m^2 - x(1-x)k^2}{\mu^2}\right)\right]
  + \text{finite},
  \label{eq:Pi_scalar_exact}
\end{equation}
where $d = 4 - \epsilon$ and $\mu$ is the renormalization scale.

In the massless limit $m \to 0$ and after renormalization (MS scheme), the log-running piece is:
\begin{equation}
  \Pi^{(1)}_{\mathrm{run}}(k^2)
  = -\frac{e^2}{6\pi^2}\,\ln\!\left(\frac{-k^2}{\mu^2}\right)\cdot\int_0^1 dx\;x(1-x)
  \;\cdot 2\pi^2 \cdot \frac{1}{2\pi^2}
  = -\frac{e^2}{6\pi^2}\,\ln\!\left(\frac{-k^2}{\mu^2}\right).
  \label{eq:Pi_scalar_run}
\end{equation}

\begin{remark}[Derivation of the $1/6$ coefficient]
The standard loop integral for a complex scalar is:
\begin{equation}
  \Pi^{(1)}(k^2) = -\frac{e^2}{2\pi^2}\int_0^1 dx\;x(1-x)\,
  \ln\!\left(\frac{m^2 - x(1-x)k^2}{\mu^2}\right) + \text{const.}
\end{equation}
Using $\int_0^1 dx\,x(1-x) = 1/6$, the coefficient of $\ln(-k^2/\mu^2)$ is:
\begin{equation}
  \frac{e^2}{2\pi^2} \cdot \frac{1}{6} = \frac{e^2}{12\pi^2}.
\end{equation}
The contribution to the running of $\alpha^{-1} = 4\pi/e^2$ is therefore:
\begin{equation}
  \frac{d\,\alpha^{-1}}{d\ln\mu}
  = -4\pi \cdot \frac{e^2}{12\pi^2} \cdot \frac{1}{\alpha^2}\cdot\alpha^2
  = -\frac{1}{3\pi}.
\end{equation}
In the UBT convention $1/\alpha(\mu) = 1/\alpha_0 + (B/2\pi)\ln(\mu/\mu_0)$, each complex
scalar contributes $\delta B = 2\pi/3$ to $B$.
\end{remark}

\subsection{$N_{\mathrm{charged}}$ Independent Scalars: Total $\Pi(k^2)$}

Since $S_{\mathrm{kin}}$ is a sum of $N_{\mathrm{charged}}$ independent kinetic terms
(Eq.~\eqref{eq:Skin_sum}), the total one-loop vacuum polarization is:
\begin{equation}
  \Pi_{\mathrm{UBT}}(k^2)
  = N_{\mathrm{charged}} \cdot \Pi^{(1)}(k^2)
  = -N_{\mathrm{charged}} \cdot \frac{e^2}{12\pi^2}\,\ln\!\left(\frac{-k^2}{\mu^2}\right)
    + \text{const.}
  \label{eq:Pi_UBT}
\end{equation}

The electromagnetic coupling running is:
\begin{equation}
  \frac{1}{\alpha(\mu)} = \frac{1}{\alpha(\mu_0)}
    + \frac{N_{\mathrm{charged}}}{3\pi}\,\ln\!\left(\frac{\mu}{\mu_0}\right).
  \label{eq:alpha_running}
\end{equation}

\subsection{Identification of $B_0$ and $N_{\mathrm{eff}}$}

The UBT effective-potential convention writes the running as:
\begin{equation}
  \frac{1}{\alpha(\mu)} = \frac{1}{\alpha(\mu_0)}
    + \frac{B_0}{2\pi}\,\ln\!\left(\frac{\mu}{\mu_0}\right).
  \label{eq:ubt_running}
\end{equation}

Comparing Eq.~\eqref{eq:alpha_running} with Eq.~\eqref{eq:ubt_running}:
\begin{equation}
  \frac{B_0}{2\pi} = \frac{N_{\mathrm{charged}}}{3\pi}
  \quad\Longrightarrow\quad
  \boxed{B_0 = \frac{2\pi\,N_{\mathrm{charged}}}{3}.}
  \label{eq:B0_Ncharged}
\end{equation}

\begin{theorem}[Identification $N_{\mathrm{eff}} = N_{\mathrm{charged}}$]
\label{thm:Neff_equals_Ncharged}
The effective mode number $N_{\mathrm{eff}}$ in the formula $B_0 = 2\pi N_{\mathrm{eff}}/3$
is precisely the number of independent charged complex degrees of freedom of $\Theta$:
\begin{equation}
  N_{\mathrm{eff}} = N_{\mathrm{charged}}.
\end{equation}
\end{theorem}

\begin{proof}
Direct comparison of Eq.~\eqref{eq:alpha_running} and Eq.~\eqref{eq:ubt_running}.
The proportionality factor between the two expressions for $d\alpha^{-1}/d\ln\mu$ is
$N_{\mathrm{charged}}/(3\pi) = B_0/(2\pi)$, which uniquely determines $B_0 = 2\pi N_{\mathrm{charged}}/3$.
\end{proof}

\subsection{QED Limit Check}

Setting $N_{\mathrm{charged}} = 1$ (single complex scalar):
\begin{equation}
  B_0^{\mathrm{QED}} = \frac{2\pi \cdot 1}{3} = \frac{2\pi}{3}, \qquad \checkmark
  \label{eq:qed_limit_B0}
\end{equation}
which reproduces the standard QED one-loop coefficient $\beta_1 = 2/3$ (since
$B_0/(2\pi) = 1/3 = \beta_1/2$). \checkmark

\subsection{Summary of Step 2}

\begin{theorem}[$B_0$ from $S_{\mathrm{kin}}$]
\label{thm:B0_from_Skin}
From the kinetic action $S_{\mathrm{kin}}[\Theta, A]$ with $\Theta \in \C \otimes \H$:
\begin{enumerate}
  \item[(i)] The one-loop vacuum polarization is
    $\Pi_{\mathrm{UBT}}(k^2) = -N_{\mathrm{charged}} \cdot \frac{e^2}{12\pi^2}\ln(-k^2/\mu^2)$.
  \item[(ii)] The $B$-coefficient is $B_0 = 2\pi N_{\mathrm{eff}}/3$ with
    $N_{\mathrm{eff}} = N_{\mathrm{charged}}$.
  \item[(iii)] QED limit: $N_{\mathrm{charged}} = 1 \Rightarrow B_0 = 2\pi/3$. \checkmark
\end{enumerate}
\end{theorem}

\medskip\noindent
\textbf{Key conclusion:} The question of the value of $N_{\mathrm{eff}}$ reduces entirely to
counting $N_{\mathrm{charged}}$ — the number of independent charged complex d.o.f.\ in $\Theta$.
This counting is performed in Step~3 (see \texttt{step3\_counting.tex}).

\ifdefined\INCLUDEMODE\else
\end{document}
\fi
